\subsection{Quotient Spaces}
\begin{definition} \label{def:affine-subset-parallel-subspace} \hlt{Affine Subset and Parallel Subspaces}\\
Let $\mathcal{U}$ be a subspace of $\mathcal{V}$. For $\mathbf{v}\in\mathcal{V}$, the subset $\mathbf{v}+\mathcal{U}\defeq\{\mathbf{v}+\mathbf{u}:\mathbf{u}\in\mathcal{U}\}$ is an \hlt{affine subset} of $\mathcal{V}$ parallel to $\mathcal{U}$.
\end{definition}

\begin{proposition} \label{prop:parallel-affine-subsets-equal-or-disjoint} \hlt{Parallel Affine Subsets Are Equal or Disjoint}\\
Let $\mathcal{U}$ be a subspace of $\mathcal{V}$ and let $\mathbf{v},\mathbf{w}\in\mathcal{V}$. The following are equivalent:
\begin{enumerate}[label=(\roman*)]
\setlength{\itemsep}{0pt}
\item $\mathbf{v}-\mathbf{w}\in\mathcal{U}$;
\item $\mathbf{v}+\mathcal{U}=\mathbf{w}+\mathcal{U}$;
\item $(\mathbf{v}+\mathcal{U})\cap(\mathbf{w}+\mathcal{U})\neq\emptyset$.
\end{enumerate}
\begin{proof}
\begin{enumerate}[label=(\roman*)]
\setlength{\itemsep}{0pt}
\item[$(i)$]$\Rightarrow (ii)$ If $\mathbf{v}-\mathbf{w}\in\mathcal{U}$ and $\mathbf{u}\in\mathcal{U}$, then $\mathbf{v}+\mathbf{u}=\mathbf{w}+((\mathbf{v}-\mathbf{w})+\mathbf{u})\in\mathbf{w}+\mathcal{U}$. Thus $\mathbf{v}+\mathcal{U}\subseteq\mathbf{w}+\mathcal{U}$, and the reverse inclusion is similar.
\item[$(ii)$]$\Rightarrow (iii)$ Equal affine subsets have the same elements.
\item[$(iii)$]$\Rightarrow (i)$ If $\mathbf{x}\in(\mathbf{v}+\mathcal{U})\cap(\mathbf{w}+\mathcal{U})$, then $\mathbf{x}=\mathbf{v}+\mathbf{u}_1=\mathbf{w}+\mathbf{u}_2$ for some $\mathbf{u}_1,\mathbf{u}_2\in\mathcal{U}$. Hence $\mathbf{v}-\mathbf{w}=\mathbf{u}_2-\mathbf{u}_1\in\mathcal{U}$.
\end{enumerate}
\hlt{Structural Commentary}: Parallel affine subsets either coincide or do not meet.\\
\hlt{Key Insight}: Intersection is controlled by the difference of base points.
\end{proof}
\end{proposition}

\begin{definition} \label{def:quotient-space} \hlt{Quotient Space}\\
Let $\mathcal{U}$ be a subspace of $\mathcal{V}$.\\
The \hlt{quotient space} $\mathcal{V}/\mathcal{U}$ is the set of all affine subsets of $\mathcal{V}$ parallel to $\mathcal{U}$: $\mathcal{V}/\mathcal{U}\defeq\{\mathbf{v}+\mathcal{U}:\mathbf{v}\in\mathcal{V}\}$.
\end{definition}

\begin{proposition} \label{prop:quotient-relation-congruence} \hlt{Quotient Relation Is a Linear Congruence}\\
Let $\mathcal{U}\leq\mathcal{V}$ and define $\mathbf{v}\equiv_{\mathcal{U}}\mathbf{w}\Longleftrightarrow \mathbf{v}-\mathbf{w}\in\mathcal{U}$.\\
Then $\equiv_{\mathcal{U}}$ is an equivalence relation and, for $a,b\in\mathbb{F}$, $\mathbf{v}_1\equiv_{\mathcal{U}}\mathbf{v}_2$, $\mathbf{w}_1\equiv_{\mathcal{U}}\mathbf{w}_2$ $\Longrightarrow$ $a\mathbf{v}_1+b\mathbf{w}_1\equiv_{\mathcal{U}} a\mathbf{v}_2+b\mathbf{w}_2$.
\begin{proof}
Since $\mathcal{U}\leq\mathcal{V}$, $\mathbf{0}\in\mathcal{U}$, $-\mathcal{U}=\mathcal{U}$, and $\mathcal{U}+\mathcal{U}\subseteq\mathcal{U}$; hence $\equiv_{\mathcal{U}}$ is reflexive, symmetric, and transitive.\\
If $\mathbf{v}_1\equiv_{\mathcal{U}}\mathbf{v}_2$ and $\mathbf{w}_1\equiv_{\mathcal{U}}\mathbf{w}_2$, then $(a\mathbf{v}_1+b\mathbf{w}_1)-(a\mathbf{v}_2+b\mathbf{w}_2) =a(\mathbf{v}_1-\mathbf{v}_2)+b(\mathbf{w}_1-\mathbf{w}_2)\in\mathcal{U}$.\\
\hlt{Structural Commentary}: Subspace congruence makes quotient operations well-defined.\\
\hlt{Key Insight}: Subspace closure is exactly the well-definedness mechanism for quotient operations.
\end{proof}
\end{proposition}

\begin{perspective} \label{per:quotient-cosets-as-equivalence-classes} \hlt{Cosets as Equivalence Classes}\\
The quotient relation is $\mathbf{v}\equiv_{\mathcal{U}}\mathbf{w}$ when $\mathbf{v}-\mathbf{w}\in\mathcal{U}$. Proposition \ref{prop:parallel-affine-subsets-equal-or-disjoint} shows that the equivalence class of $\mathbf{v}$ is exactly the coset $\mathbf{v}+\mathcal{U}$. Thus the quotient space remembers vectors only up to directions in $\mathcal{U}$.
\end{perspective}

\begin{proposition} \label{prop:quotient-space-structure} \hlt{Structure of the Quotient}\\
Let $\mathcal{U}$ be a subspace of $\mathcal{V}$. Addition and scalar multiplication on $\mathcal{V}/\mathcal{U}$ defined by $(\mathbf{v} + \mathcal{U}) + (\mathbf{w} + \mathcal{U}) \defeq (\mathbf{v} + \mathbf{w}) + \mathcal{U}$ and $a(\mathbf{v} + \mathcal{U}) \defeq a\mathbf{v} + \mathcal{U}$ are well-defined, and with these operations $\mathcal{V}/\mathcal{U}$ is a vector space.
\begin{proof}
If $\mathbf{v} + \mathcal{U} = \mathbf{v}' + \mathcal{U}$ and $\mathbf{w} + \mathcal{U} = \mathbf{w}' + \mathcal{U}$, then $\mathbf{v}-\mathbf{v}' \in \mathcal{U}$ and $\mathbf{w}-\mathbf{w}' \in \mathcal{U}$.\\
Since $\mathcal{U}$ is a subspace, $(\mathbf{v}+\mathbf{w}) - (\mathbf{v}'+\mathbf{w}') \in \mathcal{U}$, and $a\mathbf{v} - a\mathbf{v}' \in \mathcal{U}$ so the operations do not depend on the chosen representatives. The vector space axioms are inherited from $\mathcal{V}$ under these operations.\\
\hlt{Structural Commentary}: Quotients preserve operations modulo $\mathcal{U}$.\\
\hlt{Key Insight}: Cosets identify vectors differing by elements of $\mathcal{U}$.
\end{proof}
\end{proposition}

\begin{definition} \hlt{Quotient Map}\\
Let $\mathcal{U}$ be a subspace of $\mathcal{V}$. The \hlt{quotient map} is the map $\pi : \mathcal{V} \to \mathcal{V}/\mathcal{U}$ defined by $\pi(\mathbf{v})=\mathbf{v}+\mathcal{U}$.
\end{definition}

\begin{proposition} \label{prop:quotient-map-properties} \hlt{Properties of the Quotient Map}\\
The quotient map $\pi$ is linear, surjective, and $\ker(\pi) = \mathcal{U}$.
\begin{proof}
For $\mathbf{v},\mathbf{w}\in\mathcal{V}$ and $a\in\mathbb{F}$, $\pi(a\mathbf{v}+\mathbf{w})=(a\mathbf{v}+\mathbf{w})+\mathcal{U} =a(\mathbf{v}+\mathcal{U})+(\mathbf{w}+\mathcal{U}) =a\pi(\mathbf{v})+\pi(\mathbf{w})$ so $\pi$ is linear. It is surjective because every coset in $\mathcal{V}/\mathcal{U}$ has the form $\mathbf{v}+\mathcal{U}=\pi(\mathbf{v})$. Finally, $\ker(\pi)=\{\mathbf{v}\in\mathcal{V}:\mathbf{v}+\mathcal{U}=\mathcal{U}\}=\mathcal{U}$.\\
\hlt{Structural Commentary}: The quotient map collapses $\mathcal{U}$ to zero and records cosets.\\
\hlt{Key Insight}: Quotient maps realise subspaces as kernels.
\end{proof}
\end{proposition}

\begin{theorem} \label{thm:dimension-of-quotient} \hlt{Dimension of a Quotient}\\
If $\mathcal{V}$ is finite-dimensional and $\mathcal{U}$ is a subspace, then $\dim(\mathcal{V}/\mathcal{U}) = \dim(\mathcal{V}) - \dim(\mathcal{U})$.
\begin{proof}
Choose a basis $(\mathbf{u}_1,\ldots,\mathbf{u}_m)$ of $\mathcal{U}$ and extend it to a basis $(\mathbf{u}_1,\ldots,\mathbf{u}_m,\mathbf{w}_1,\ldots,\mathbf{w}_k)$ of $\mathcal{V}$.\\
The cosets $\mathbf{w}_1+\mathcal{U},\ldots,\mathbf{w}_k+\mathcal{U}$ span $\mathcal{V}/\mathcal{U}$ because the $\mathbf{u}_j$-terms vanish in the quotient.\\
If $a_1(\mathbf{w}_1+\mathcal{U})+\cdots+a_k(\mathbf{w}_k+\mathcal{U})=\mathcal{U}$, then $a_1\mathbf{w}_1+\cdots+a_k\mathbf{w}_k\in\mathcal{U}$, so a linear relation in the extended basis forces $a_1=\cdots=a_k=0$. Hence these cosets form a basis of $\mathcal{V}/\mathcal{U}$, and its length is $k=\dim(\mathcal{V})-\dim(\mathcal{U})$.\\
\hlt{Structural Commentary}: Quotiening removes the directions in $\mathcal{U}$ and keeps the added basis directions.\\
\hlt{Key Insight}: Passing to $\mathcal{V}/\mathcal{U}$ subtracts the collapsed dimension.
\end{proof}
\end{theorem}

\begin{theorem} \label{thm:quotient-universal-property} \label{thm:quotient-factorisation-through-projection} \hlt{Universal Property of the Quotient Map}\\
Let $\mathcal{U}\leq\mathcal{V}$, $\pi:\mathcal{V}\to\mathcal{V}/\mathcal{U}$, and $T\in\LL(\mathcal{V},\mathcal{W})$ with $\mathcal{U}\leq\ker(T)$.\\
Then $\exists!\,\overline{T}\in\LL(\mathcal{V}/\mathcal{U},\mathcal{W})
$ such that $\overline{T}\pi=T$, $\overline{T}(\mathbf{v}+\mathcal{U})=T\mathbf{v}$, and $\ker(\overline{T})=\ker(T)/\mathcal{U}$, $\operatorname{range}(\overline{T})=\operatorname{range}(T)$.
\begin{proof}
Define $\overline{T}(\mathbf{v}+\mathcal{U})=T\mathbf{v}$. If $\mathbf{v}+\mathcal{U}=\mathbf{w}+\mathcal{U}$, then $\mathbf{v}-\mathbf{w}\in\mathcal{U}\leq\ker(T)$, so $T\mathbf{v}=T\mathbf{w}$; hence $\overline{T}$ is well-defined.\\
Also $\overline{T}\big(a(\mathbf{v}+\mathcal{U})+b(\mathbf{w}+\mathcal{U})\big) =T(a\mathbf{v}+b\mathbf{w}) =aT\mathbf{v}+bT\mathbf{w}$.\\
By definition, $\overline{T}\pi=T$. If $S\pi=T$, then $S(\mathbf{v}+\mathcal{U})=S\pi(\mathbf{v})=T\mathbf{v}=\overline{T}(\mathbf{v}+\mathcal{U})$, so $S=\overline{T}$.\\
Also $\mathbf{v}+\mathcal{U}\in\ker(\overline{T}) \Longleftrightarrow T\mathbf{v}=\mathbf{0} \Longleftrightarrow \mathbf{v}\in\ker(T)$.\\
So $\ker(\overline{T})=\ker(T)/\mathcal{U}$, and $\operatorname{range}(\overline{T})=\{\overline{T}(\mathbf{v}+\mathcal{U}):\mathbf{v}\in\mathcal{V}\}=\operatorname{range}(T)$.\\
\hlt{Structural Commentary}: Maps killing $\mathcal{U}$ factor through $\mathcal{V}/\mathcal{U}$.\\
\hlt{Key Insight}: Quotients are universal domains for maps with prescribed kernel ambiguity.
\end{proof}
\end{theorem}

\begin{figure}[H]
  \centering
  \begin{tikzpicture}[
  x=0.88cm,
  y=0.88cm,
  >=Stealth,
  every node/.style={font=\small},
  point/.style={circle, fill=gray!75, inner sep=1.55pt},
  arrow/.style={->, gray!65!black, line width=0.7pt},
  factor/.style={->, dashed, gray!65!black, line width=0.7pt},
  hint/.style={font=\scriptsize, text=gray!70!black, inner sep=1pt}
]
  \node[point] (V) at (0,0.75) {};
  \node[point] (Q) at (0,-0.75) {};
  \node[point] (W) at (3.1,0.75) {};

  \draw[arrow] (V) -- node[above, hint] {$L$} (W);
  \draw[arrow] (V) -- node[left, hint] {$\pi$} (Q);
  \draw[factor] (Q) -- node[above, sloped, hint] {$\overline{L}$} (W);

  \node at (0,1.08) {$\mathcal{V}$};
  \node at (3.1,1.08) {$\mathcal{W}$};
  \node at (0,-1.1) {$\mathcal{V}/\mathcal{U}$};
\end{tikzpicture}

  \caption{If $\mathcal{U}\subseteq\ker(L)$, then $L=\overline{L}\pi$ uniquely.}
  \label{fig:quotient-universal-property}
\end{figure}

\begin{theorem} \label{thm:subspace-correspondence-quotient} \hlt{Subspace Correspondence for Quotients}\\
Let $\pi:\mathcal{V}\to\mathcal{V}/\mathcal{U}$ be the quotient map. The assignments $\mathcal{T}\longmapsto \mathcal{T}/\mathcal{U}$ and $\mathcal{Y}\longmapsto \pi^{-1}(\mathcal{Y})$ are inverse inclusion-preserving bijections between $\{\mathcal{T}$ subspace of $\mathcal{V}:\mathcal{U}\subseteq\mathcal{T}\}$ and $\{\mathcal{Y}$ subspace of $\mathcal{V}/\mathcal{U}\}$.
\begin{proof}
If $\mathcal{U}\subseteq\mathcal{T}\subseteq\mathcal{V}$, then $\mathcal{T}/\mathcal{U}$ is a subspace of $\mathcal{V}/\mathcal{U}$.\\
If $\mathcal{Y}\subseteq\mathcal{V}/\mathcal{U}$ is a subspace, then $\pi^{-1}(\mathcal{Y})$ is a subspace of $\mathcal{V}$ and contains $\ker(\pi)=\mathcal{U}$.\\
For $\mathcal{U}\subseteq\mathcal{T}\subseteq\mathcal{V}$, $\pi^{-1}(\mathcal{T}/\mathcal{U})=\mathcal{T}$, since $\mathbf{v}+\mathcal{U}=\mathbf{t}+\mathcal{U}$ with $\mathbf{t}\in\mathcal{T}$ implies $\mathbf{v}-\mathbf{t}\in\mathcal{U}\subseteq\mathcal{T}$.\\
For $\mathcal{Y}\subseteq\mathcal{V}/\mathcal{U}$, surjectivity gives $\pi\left(\pi^{-1}(\mathcal{Y})\right)=\mathcal{Y}$, that is, $\pi^{-1}(\mathcal{Y})/\mathcal{U}=\mathcal{Y}$\\
Both assignments preserve inclusion by monotonicity of images and preimages.\\
\hlt{Structural Commentary}: Subspaces above $\mathcal{U}$ are exactly the subspaces visible after collapsing $\mathcal{U}$.\\
\hlt{Key Insight}: The quotient map translates containment above $\mathcal{U}$ into containment inside $\mathcal{V}/\mathcal{U}$.
\end{proof}
\end{theorem}

\begin{figure}[H]
  \centering
  \begin{tikzpicture}[
  x=0.88cm,
  y=0.88cm,
  >=Stealth,
  every node/.style={font=\small},
  point/.style={circle, fill=gray!75, inner sep=1.55pt},
  edge/.style={->, gray!65!black, line width=0.7pt},
  map/.style={->, gray!60, line width=0.65pt},
  hint/.style={font=\scriptsize, text=gray!70!black, inner sep=1pt}
]
  \node[point] (U) at (-2.4,0.55) {};
  \node[point] (T) at (0,0.55) {};
  \node[point] (V) at (2.4,0.55) {};

  \node[point] (QU) at (-2.4,-0.55) {};
  \node[point] (QT) at (0,-0.55) {};
  \node[point] (QV) at (2.4,-0.55) {};

  \draw[edge] (U) -- (T);
  \draw[edge] (T) -- (V);
  \draw[edge] (QU) -- (QT);
  \draw[edge] (QT) -- (QV);

  \draw[map] (U) -- (QU);
  \draw[map] (T) -- node[right, hint] {$\pi$} (QT);
  \draw[map] (V) -- (QV);

  \node at (-2.4,0.9) {$\mathcal{U}$};
  \node at (0,0.9) {$\mathcal{T}$};
  \node at (2.4,0.9) {$\mathcal{V}$};

  \node at (-2.4,-0.92) {$\{\mathcal{U}\}$};
  \node at (0,-0.92) {$\mathcal{T}/\mathcal{U}$};
  \node at (2.4,-0.92) {$\mathcal{V}/\mathcal{U}$};
\end{tikzpicture}

  \caption{Intermediate subspaces above $\mathcal{U}$ correspond to subspaces of the quotient.}
  \label{fig:quotient-correspondence-subspaces}
\end{figure}

\begin{theorem} \label{thm:first-isomorphism-theorem-vector-spaces} \label{thm:quotient-by-kernel-is-range} \hlt{Quotient by the Kernel}\\
Let $T\in\LL(\mathcal{V},\mathcal{W})$. Then $\widetilde{T}(\mathbf{v}+\ker(T))\defeq T\mathbf{v}$ determines a well-defined linear map $\widetilde{T}:\mathcal{V}/\ker(T)\to\mathcal{W}$. Also,
\begin{enumerate}[label=(\roman*)]
\setlength{\itemsep}{0pt}
\item $\widetilde{T}$ is injective;
\item $\operatorname{range}(\widetilde{T})=\operatorname{range}(T)$;
\item $\mathcal{V}/\ker(T)\cong\operatorname{range}(T)$.
\end{enumerate}
\begin{proof}
If $\mathbf{v}+\ker(T)=\mathbf{v}'+\ker(T)$, then $\mathbf{v}-\mathbf{v}'\in\ker(T)$, so $T\mathbf{v}=T\mathbf{v}'$. Thus $\widetilde{T}$ is well-defined.\\
For cosets and $a\in\mathbb{F}$, $\widetilde{T}\big(a(\mathbf{v}+\ker(T))+(\mathbf{v}'+\ker(T))\big) =T(a\mathbf{v}+\mathbf{v}') =aT\mathbf{v}+T\mathbf{v}'$, so $\widetilde{T}$ is linear.
\begin{enumerate}[label=(\roman*)]
\setlength{\itemsep}{0pt}
\item If $\widetilde{T}(\mathbf{v}+\ker(T))=\mathbf{0}$, then $T\mathbf{v}=\mathbf{0}$, so $\mathbf{v}\in\ker(T)$ and $\mathbf{v}+\ker(T)=\ker(T)$, the zero coset.\\
Hence $\ker(\widetilde{T})=\{\ker(T)\}$ and $\widetilde{T}$ is injective.
\item For every coset $\mathbf{v}+\ker(T)$, the value $\widetilde{T}(\mathbf{v}+\ker(T))=T\mathbf{v}$ lies in $\operatorname{range}(T)$.\\
Conversely, every vector in $\operatorname{range}(T)$ has the form $T\mathbf{v}=\widetilde{T}(\mathbf{v}+\ker(T))$. Hence the ranges are equal.
\item Let $\widehat{T}:\mathcal{V}/\ker(T)\to\operatorname{range}(T)$ be the same formula as $\widetilde{T}$ with codomain restricted to $\operatorname{range}(T)$.\\
The linearity proved above and parts (i)--(ii) show that $\widehat{T}$ is linear, injective, and surjective, so Proposition \ref{prop:linear-map-invertible-iff-bijective} gives $\mathcal{V}/\ker(T)\cong\operatorname{range}(T)$.
\end{enumerate}
\hlt{Structural Commentary}: Quotienting by the kernel removes exactly the ambiguity of $T$.\\
\hlt{Key Insight}: Every linear map factors as a quotient followed by an isomorphism onto its range.
\end{proof}
\end{theorem}

\begin{theorem} \label{thm:second-isomorphism-vector-spaces} \hlt{Second Isomorphism Theorem}\\
Let $\mathcal{S}$ and $\mathcal{T}$ be subspaces of $\mathcal{V}$. Then $(\mathcal{S}+\mathcal{T})/\mathcal{T}\cong \mathcal{S}/(\mathcal{S}\cap\mathcal{T})$.
\begin{proof}
Define $L:\mathcal{S}\to(\mathcal{S}+\mathcal{T})/\mathcal{T}$ by $L\mathbf{s}=\mathbf{s}+\mathcal{T}$. Then $L$ is linear by the quotient operations.\\
Moreover, $\ker(L)=\{\mathbf{s}\in\mathcal{S}:\mathbf{s}+\mathcal{T}=\mathcal{T}\} =\mathcal{S}\cap\mathcal{T}$.
If $(\mathbf{s}+\mathbf{t})+\mathcal{T}\in(\mathcal{S}+\mathcal{T})/\mathcal{T}$, with $\mathbf{s}\in\mathcal{S}$ and $\mathbf{t}\in\mathcal{T}$, then $(\mathbf{s}+\mathbf{t})+\mathcal{T}=\mathbf{s}+\mathcal{T}=L\mathbf{s}$, so $L$ is surjective. Theorem \ref{thm:first-isomorphism-theorem-vector-spaces} gives the stated isomorphism.\\
\hlt{Structural Commentary}: Modding out by $\mathcal{T}$ leaves only the part of $\mathcal{S}$ not already inside $\mathcal{T}$.\\
\hlt{Key Insight}: In a quotient by $\mathcal{T}$, adding $\mathcal{T}$ to $\mathcal{S}$ costs nothing.
\end{proof}
\end{theorem}

\begin{theorem} \label{thm:third-isomorphism-vector-spaces} \hlt{Third Isomorphism Theorem}\\
If $\mathcal{U}\subseteq\mathcal{T}\subseteq\mathcal{V}$ are subspaces, then $(\mathcal{V}/\mathcal{U})/(\mathcal{T}/\mathcal{U})\cong \mathcal{V}/\mathcal{T}$.
\begin{proof}
Define $L:\mathcal{V}/\mathcal{U}\to\mathcal{V}/\mathcal{T}$ by $L(\mathbf{v}+\mathcal{U})=\mathbf{v}+\mathcal{T}$.\\
If $\mathbf{v}+\mathcal{U}=\mathbf{w}+\mathcal{U}$, then $\mathbf{v}-\mathbf{w}\in\mathcal{U}\subseteq\mathcal{T}$, so $L$ is well-defined.\\
It is linear by the quotient operations, and it is surjective since $L(\mathbf{v}+\mathcal{U})=\mathbf{v}+\mathcal{T}$ for every coset $\mathbf{v}+\mathcal{T}$.\\
Moreover, $\ker(L)=\{\mathbf{v}+\mathcal{U}: \mathbf{v}+\mathcal{T}=\mathcal{T}\} =\{\mathbf{v}+\mathcal{U}: \mathbf{v}\in\mathcal{T}\} =\mathcal{T}/\mathcal{U}$. Theorem \ref{thm:first-isomorphism-theorem-vector-spaces} gives $(\mathcal{V}/\mathcal{U})/(\mathcal{T}/\mathcal{U})\cong\mathcal{V}/\mathcal{T}$.\\
\hlt{Structural Commentary}: Collapsing $\mathcal{U}$ and then collapsing $\mathcal{T}/\mathcal{U}$ is the same as collapsing $\mathcal{T}$ at once.\\
\hlt{Key Insight}: Quotients of quotients cancel only when the nesting is explicit.
\end{proof}
\end{theorem}

This block develops two complementary ways to organise a vector space: build from smaller pieces by spans and direct sums, or collapse a chosen subspace by passing to a quotient. Bases and dimension measure exactly what is added, split off, or made invisible.
